\documentclass[11pt]{article}

% --- يجب التصريف باستخدام XeLaTeX (RTL عبر polyglossia/bidi، وخط Amiri) ---
\usepackage{amsmath,amssymb,amsthm}
\usepackage{booktabs}
\usepackage{geometry}
\usepackage{xcolor}
\usepackage{hyperref}
\geometry{margin=1in}
\hypersetup{colorlinks=true,linkcolor=blue,citecolor=blue,urlcolor=blue}

\usepackage{polyglossia}
\setmainlanguage[numerals=maghrib]{arabic}
\setotherlanguage{english}
\newfontfamily\arabicfont[Script=Arabic]{Amiri}

\newtheorem{definition}{تعريف}
\newtheorem{proposition}{قضية}
\newtheorem{theorem}{نظرية}

\newcommand{\dcausal}{d_{\mathrm{c}}}
\newcommand{\Nk}{\mathcal{N}_k}
% يُجبَر الرمز البرمجي على الاتجاه من اليسار إلى اليمين داخل النص العربي.
\newcommand{\code}[1]{\textenglish{\texttt{#1}}}

\title{\textbf{وحدة إدارة الذاكرة الإدراكية:\\ ذاكرة عمل حتمية وقابلة للتدقيق لوكلاء البرمجة المعتمدين على النماذج اللغوية الكبيرة\\ \large (وسردٌ أمينٌ لسبب عدم تفوّق الاسترجاع السببي على RAG)}}
\author{مشروع CCOS\thanks{Causal Context Operating System (CCOS)، نموذج بحثي أوّلي
مفتوح. الشيفرة المصدرية وبرامج إعادة الإنتاج وهذه الورقة:
\url{https://github.com/CHECKUPAUTO/CCOS}.}}
\date{\today}

\begin{document}
\maketitle

\begin{abstract}
تُدار الذاكرة العاملة لوكيل برمجة معتمد على نموذج لغوي كبير بطريقة مُعتِمة: إذ يقوم
مكدّس استرجاع (RAG، أو GraphRAG، أو أطر الوكلاء) بترحيل مستودعٍ إلى نافذة محدودة، ثم
يُحوّر تلك النافذة احتماليًّا مع تقدّم الجلسة، حتى إذا انحرف الوكيل بعد عشرات الخطوات
استحال أن نقول \emph{لماذا} أو \emph{متى} فسد تمثيله للمشروع. يعامل CCOS ذاكرة الوكيل
كما يعامل نظام التشغيل الذاكرة العشوائية (RAM): تُحلَّل الشيفرة المصدرية إلى رسم بياني
سببي \emph{قائم على مصدر الأحداث}، وتُسجَّل العُقد وتُرحَّل ضمن ميزانية رموز، وتُدوَّن كل
انتقالة حالة في سجلٍّ مُسلسَل بالتجزئة يُعاد تشغيله بِتةً بِتة. نقدّم تعريفًا دقيقًا
وحتميًّا لـ\emph{المنطقة السببية} (مركّبة متّصلة من رسم الروابط السببية عبر الملفات)
وللمسافة السببية (أقصر مسار موزون)، ونبرهن أنّ المناطق والنافذة المُرحَّلة دالّةٌ خالصة
في الرسم البياني، تُعاد بناؤها بالضبط عند إعادة التشغيل --- وهي الخاصية التي تجعل ذاكرة
الوكيل \emph{قابلة للتدقيق} و\emph{قابلة للتنقيح بالسفر عبر الزمن}. ثم نسأل هل تجعل هذه
البنية السببية من CCOS \emph{مُسترجِعًا أفضل} من خط أساس، ونُبلِغ --- بأمانة --- أنّها لا
تفعل. فبالقياس على $70$ تثبيتًا حقيقيًّا لإصلاح علل عبر ثلاث حُزَم (crates) ناضجة،
\textbf{يتعادل الاختيار السببي مع مُسترجِع معجمي بسيط من نوع TF-IDF}، ويخسر أمامه عند
ميزانية ضيّقة؛ ففي العلل الحقيقية تتشارك ملفات الإصلاح المفردات، فيجدها التشابه المعجمي
أيضًا، كما أنّ إثباتًا للمفهوم يُمهِّد CCOS انطلاقًا من أثر انهيار يُهزَم بدوره بـ«RAG على
رسالة الخطأ» (فرسائل خطأ لغة Rust تُسمّي السبب). أمّا محاكاةٌ من دون نموذج لغوي
\emph{فتفصل} بينهما حين يُبنى السبب ليكون معجميًّا غير مشابه لعَرَضِه، ويُعيد تشغيلٌ أوّل
على الجهاز بنموذجٍ من $7$ مليارات معامل ذلك --- لكنّ هذا النظام لا يحدث على شيفرة حقيقية.
ولذا فالإسهام \emph{ليس} استرجاعًا أفضل. إنّه البنية التحتية التي يفتقر إليها مكدّس
الاسترجاع بنيويًّا: ذاكرة وكيلٍ \emph{حتمية، قابلة لإعادة التشغيل، قابلة للتدقيق} يمكن فيها
إرجاع حالة السياق الدقيقة عند أيّ خطوة وإعادة تشغيلها بمعاملات مختلفة (\emph{التنقيح بالسفر
عبر الزمن}). نُصدِر المحرّك ومنصّة التحقّق الأمينة وهذه الورقة، مع وسمٍ في كل موضع لما هو
\emph{مُبرهَن} (البنية الصورية، والحتمية، وإعادة التشغيل)، وما هو \emph{مُقاس} (تغطية
الاسترجاع --- حيث \emph{لا} يتفوّق CCOS على RAG المعجمي في العلل الحقيقية)، والمحور الوحيد
الذي \emph{يفوز} فيه CCOS: \textbf{الكفاءة}. فعند التشغيل من الطرف إلى الطرف على إصلاحات
حقيقية أحادية الملف بنموذجٍ من حجم 30B مع إعادة محاولة و«المُصرِّف في الحلقة»، يحلّ CCOS
وRAG \emph{بالتساوي} ($2/10$)، لكنّ CCOS يفعل ذلك بـ$776$ رمز سياق مقابل $5366$
للمُسترجِع المعجمي --- خفضٌ بمقدار $6.9\times$ (و$4.1$ إلى $9.1\times$ عبر $51$ إصلاحًا من
ثلاث حُزَم، من دون نموذج)، لأنّ المنطقة السببية تتوقّف عند مجموعة العمل بدل حشو
الـ top-$k$ حتى الميزانية. فـ CCOS إذن ليس مُسترجِعًا أو حلّالًا أفضل، بل أكثر
\emph{اقتصادًا} و\emph{قابلية للتدقيق} --- يُعايِر نفسه ذاتيًّا، من دون $k$ يحتاج إلى ضبط.
(على شجرة مصدره الخاصّة، يغطّي الاختيار بالمناطق جوار مهمّةٍ بنصف قطر $k$ قفزات باستدعاء
$0.97$ مقابل $0.35$ للاختيار المسطّح، وبرموزٍ أقل بنحو $48\%$؛ والمناطق متّصلة داخليًّا
بنسبة $95.5\%$.)
\end{abstract}

\section{مقدّمة}

على وكيل برمجة معتمد على نموذج لغوي كبير يعمل فوق مستودع أن يقرّر مرارًا \emph{ما الذي
يضعه في نافذة السياق}. والنمط السائد هو الاسترجاع: تُضمَّن مقاطع الشيفرة، وتُرتَّب إزاء
المهمّة، ثم تُوصَل أعلى $k$ منها حتى تنفد ميزانية الرموز~\cite{lewis2020rag}. وهذا يعامل
السياق كقائمة مرتّبة أحادية البُعد. فينتج نمطا فشل. الأوّل، \emph{التخفيف}: مع طول المهام
تتراكم في النافذة شيفرةٌ بارزةٌ عالميًّا لكنّها غير ذات صلة بالمهمّة، والنماذج تُحسن
الانتباه ضعيفًا إلى وسط السياقات الطويلة~\cite{liu2023lost}. والثاني، \emph{عدم الاتّساق}:
ليس بالضرورة أن تشكّل مقاطع أعلى $k$ وحدةً متّصلة من الاستدلال --- فقد تُرحَّل دالّةٌ من
دون مُستدعِيها، أو مواضع أخطائها، أو البيانات التي تعتمد عليها.

واجهت أنظمة التشغيل المشكلة المماثلة قبل عقود وأجابت عنها بـ\emph{مجموعة العمل}: مجموعة
الصفحات المُشار إليها حديثًا والتي يجب أن تبقى مقيمة~\cite{denning1968working}. ويتبنّى
CCOS هذا التشبيه لسياق النموذج اللغوي~\cite{packer2023memgpt}: تُحلَّل الشيفرة المصدرية إلى
رسم بياني سببي، وتُسجَّل العُقد، وتُرحَّل «نافذة سياق» محدودة دخولًا وخروجًا مثل
RAM~$\leftrightarrow$~VRAM. وخلافًا لمكدّس الاسترجاع، تُدوَّن كل انتقالة في سجلّ أحداث
مُسلسَل بالتجزئة، فلا تكون الذاكرة صندوقًا أسود بل أثرًا \emph{قابلًا للتدقيق ولإعادة
التشغيل}. وإسهاماتنا هي:

\begin{enumerate}
  \item \textbf{نموذجٌ صوريٌّ وحتميٌّ للمناطق السياقية} (\S\ref{sec:regions},
  \S\ref{sec:determinism}): المسافة السببية كأقصر مسار موزون، والانتماء إلى المنطقة
  كمركّبة متّصلة من رسم الروابط السببية عبر الملفات، ونظرية حتمية --- فالمناطق والنافذة
  المُرحَّلة دالّةٌ خالصة في الرسم، فتُعاد جلسةٌ كاملة بِتةً بِتة من سجلّ أحداثها المُسلسَل
  بالتجزئة (إعادة تشغيل تكشف العبث).
  \item \textbf{جلسات وكيلٍ قائمة على مصدر الأحداث مع التنقيح بالسفر عبر الزمن}
  (\S\ref{sec:timetravel}): تُسجَّل كل عملية إدراكية (استيعاب، إشارة فشل، استدعاء)؛
  وتُعاد بناء حالة السياق الدقيقة عند أيّ خطوة؛ ويمكن \emph{إعادة تشغيل} استدعاءٍ بمعاملات
  مختلفة للسؤال عمّا إذا كان الوكيل سيقرّر أفضل --- وهي القدرة التي يفتقر إليها مكدّس
  الاسترجاع الاحتمالي بنيويًّا.
  \item \textbf{منصّة تحقّقٍ أمينة ونتيجةٌ سلبية} (\S\ref{sec:protocol}): على $70$ تثبيتًا
  حقيقيًّا لإصلاح علل، \emph{لا} يتفوّق الاختيار السببي على مُسترجِع معجمي من نوع TF-IDF في
  وضع ملفات الإصلاح داخل النافذة (يتعادل، ويخسر عند ميزانية ضيّقة)، ويُهزَم محورُ أثر
  الانهيار بـ«RAG على رسالة الخطأ». نُبلِغ بهذا صراحةً --- إذ يَنقل قيمة CCOS من
  \emph{الاسترجاع} إلى \emph{قابلية التدقيق}.
  \item \textbf{قياسٌ للمحلّية من دون نموذج لغوي} (\S\ref{sec:eval}) بأرقامٍ حقيقية وقابلة
  للإعادة، وبروتوكول شرطٍ \emph{كافٍ} قابلٍ للدحض (المرحلة~4) نُحدِّده لكن نتركه مفتوحًا.
\end{enumerate}

\section{أعمالٌ ذات صلة}\label{sec:related}

\paragraph{التوليد المُعزَّز بالاسترجاع.} يُعزِّز RAG نموذجًا بارامتريًّا بمخزنٍ غير
بارامتري ويسترجع مقاطع لكل استعلام~\cite{lewis2020rag}. ويضيف Self-RAG رموز تأمّل تتحكّم
في الاسترجاع وتنقد التوليدات~\cite{asai2023selfrag}. وهذه تعمل على مقاطع \emph{مستقلّة}؛
ولا يُنمذَج الاتّساق بين العناصر المسترجَعة.

\paragraph{استرجاعٌ واعٍ بالرسوم والبنية.} يبني GraphRAG رسمًا معرفيًّا للكيانات ويجيب عن
الاستعلامات العامّة بتلخيص بنية المجتمعات~\cite{edge2024graphrag}. وللشيفرة، توحِّد رسوم
الخصائص شجرة AST وتدفّق التحكّم وتدفّق البيانات~\cite{yamaguchi2014cpg}. ومناطق CCOS من
هذه السلالة لكنّها تستهدف \emph{الترحيل}: أيّ بنيةٍ فرعية متّصلة تُجعَل مقيمة لمهمّةٍ ما،
ضمن ميزانية رموز، مع إخلاءٍ حتمي.

\paragraph{ذاكرة الوكيل.} يصوّر MemGPT النموذج اللغوي كنظام تشغيلٍ يُرحِّل بين نافذةٍ ضمن
السياق ومخزنٍ خارجي~\cite{packer2023memgpt}؛ ويسترجع Generative Agents ذكرياتٍ مُسجَّلة
بالحداثة والأهمّية والصلة~\cite{park2023generative} --- وهي العوامل نفسها التي يجمعها CCOS
في حرارة منطقة. ويُنظِّم LangGraph~\cite{langgraph} الوكلاء كرسوم خطواتٍ ذات حالة؛ فهو
يُنسِّق تدفّق \emph{التحكّم}، بينما يُنظِّم CCOS \emph{الذاكرة}. وكلاهما متكامل.

\paragraph{إدارة نافذة السياق.} يُبقي إخلاء مخبأ KV الرموز «الثقيلة الوقع» أو رموز
المَصبّ مقيمة~\cite{liu2023lost}، ويطبّق PagedAttention ترحيل نظام التشغيل على مخبأ
KV~\cite{kwon2023paged}. وتعمل هذه التقنيات على مستوى الرمز داخل تمريرةٍ أمامية واحدة؛
بينما يعمل CCOS على المستوى \emph{الدلالي} عبر جلسة وكيلٍ كاملة.

\section{الركيزة السياقية السببية}\label{sec:substrate}

يُحلِّل CCOS الشيفرة المصدرية إلى \emph{رسم ذاكرة سببي} موجَّه $G=(V,E)$. والعُقدة $v\in V$
ملفٌّ أو وحدةٌ نمطية أو رمزٌ أو اعتماديةٌ خارجية، ذات حقول قياسية $\mathrm{imp}(v)$
(الأهمّية الأساس) و$\mathrm{fail}(v)\in[0,1]$ (صلة الفشل) و$\mathrm{rec}(v)\in[0,1]$
(الحداثة). وتحمل الحافة $e=(u\!\to\!w)\in E$ وزنًا $w(e)\in(0,1]$ ونوعًا (احتواء،
اعتمادية، إحالة، سببية). تُسنِد النواة لكل عقدة درجةً سببية
\begin{equation}
\mathrm{score}(v) = \mathrm{clamp}\big(0.15\,\mathrm{imp}(v) + 0.50\,\mathrm{fail}(v)
  + 0.30\,\mathrm{rec}(v) + 0.05\ln(1{+}\mathrm{acc}(v)),\,0,1\big),
\label{eq:score}
\end{equation}
حيث $\mathrm{acc}(v)$ عدّاد الوصول. تنتشر الأعطال على طول الحواف، وتتلاشى $\mathrm{fail}$
مع ساعةٍ منطقية، وتُلحَق كل انتقالة حالة بسجلّ أحداثٍ مُسلسَل بالتجزئة يُتيح إعادة تشغيلٍ
حتمية~\cite{haber1991timestamp}. ومعرّفات العُقد ذات فضاء أسماء (\code{file:p}،
\code{mod:p:n}، \code{use:p:path}، \code{sym:p:n}، \code{dep:root})؛ والملفّ المالك لعقدةٍ
قابلٌ للاسترداد من معرّفها. ويُرحِّل CCOS v0.2 العُقد ذات أعلى $\mathrm{score}$: سياسةٌ
مسطّحة أحادية البُعد. ونجعل الاختيار الآن مكانيًّا.

\section{المناطق السياقية}\label{sec:regions}

\subsection{المسافة السببية}

\begin{definition}[المسافة السببية]
ليكن $\hat G$ الرسم متعدّد الحواف غير الموجَّه على $V$ المُستحَثّ بـ$E$، ولنُسنِد لكل حافة
الكلفة $c(e) = -\ln w(e) \ge 0$، بحيث تكون الرابطة السببية الأقوى خطوةً أقصر.
\emph{المسافة السببية} $\dcausal(u,v)$ هي أدنى كلفة كلّية على كل مسارات $u$--$v$ في
$\hat G$، و$+\infty$ إن لم يوجد أيّ مسار. وتُعرَّف مسافة القفزات غير الموزونة
$\mathrm{hops}(u,v)$ على نحوٍ مماثل بكلفٍ وحدوية.
\end{definition}

بما أنّ $w(e)\le 1$ فإنّ $c(e)\ge 0$، ومن ثَمّ فإنّ $\dcausal$ مقياسُ أقصر مسارٍ حقيقي على
كل مركّبة متّصلة (غير سالب، متناظر، يحقّق متباينة المثلّث). و\emph{الجوار السببي بنصف قطر
$k$ قفزات} لهدفٍ $t$ هو
\begin{equation}
\Nk(t) = \{\, v \in V : \mathrm{hops}(t,v) \le k \,\}.
\end{equation}
و$\Nk(t)$ هو الحقيقة المرجعية «ما هو ذو صلةٍ سببيًّا بمهمّةٍ عند $t$» المستخدَمة في التقييم
(\S\ref{sec:eval}).

\subsection{الانتماء إلى المنطقة}

تربط محاور الاعتماديات الخارجية (مثل \code{dep:std}) كلّ شيءٍ تقريبًا ويجب ألّا تُطبِق
الرسم في مركّبةٍ واحدة. ولذا نفصلها.

\begin{definition}[رابطة سببية عبر الملفات]
لعقدةٍ غير خارجية $v$ ليكن $\phi(v)$ ملفَّها المالك. يكون الملفّان $f,g$ \emph{مرتبطين
مباشرة}، ويُكتَب $f \approx g$، إذا وفقط إذا وُجدت حافةٌ $(u\!\to\!w)\in E$ بحيث
$\phi(u)=f$ و$\phi(w)=g$ و$f\neq g$، ولم يكن أيٌّ من $u$ أو $w$ عقدةَ اعتماديةٍ خارجية.
\end{definition}

\begin{definition}[المنطقة]\label{def:region}
ليكن $\approx^{*}$ الإغلاق الانعكاسي المتعدّي لـ$\approx$ على مجموعة مفاتيح الملفات.
و\emph{المنطقة} هي مجموعة كل العُقد التي تقع ملفّاتها في صنف تكافؤٍ واحد من $\approx^{*}$؛
وتشكّل كل عُقد الاعتماديات الخارجية منطقةً إضافية واحدة. وتنتمي عقدتان إلى المنطقة نفسها
إذا وفقط إذا كانت ملفّاتهما في المركّبة المتّصلة نفسها من رسم الروابط السببية عبر الملفات.
\end{definition}

\begin{proposition}[المناطق تُجزّئ العُقد]
تشكّل مناطق التعريف~\ref{def:region} تجزئةً لـ$V$: كل عقدةٍ تقع في منطقةٍ واحدة بالضبط.
\end{proposition}
\begin{proof}
$\approx^{*}$ علاقة تكافؤ على مفاتيح الملفات، فأصنافها تُجزّئ المفاتيح؛ والدلو الخارجي صنفٌ
متمايز بالبناء. وربط كل عقدةٍ غير خارجية بصنف ملفّها وكل عقدةٍ خارجية بالصنف الخارجي دالّةٌ
كلّية نحو كُتلٍ منفصلة، ومن ثَمّ تجزئة. \qed
\end{proof}

افتراضيًّا (حواف احتواءٍ واستيرادٍ فقط) يكون كل ملفّ مصدري منطقته الخاصّة. أمّا اعتماديةٌ
حقيقية عبر الملفات أو فشلٌ منتشر فيُدمِج عدّة ملفات في منطقةٍ واحدة متعدّدة الملفات --- وهي
«منطقة المعرفة» التي ينبغي للوكيل أن يوقظها معًا.

\subsection{قياسات المنطقة}

لمنطقةٍ $R$ ذات مجموعة أعضاء $M$:
\begin{align}
\mathrm{heat}(v)        &= 0.5\,\mathrm{score}(v) + 0.3\,\mathrm{fail}(v) + 0.2\,\mathrm{rec}(v), \\
\mathrm{temp}(R)        &= \mathrm{clamp}\!\Big(\tfrac{1}{|M|}\textstyle\sum_{v\in M}\mathrm{heat}(v),\,0,1\Big), \label{eq:temp}\\
\mathrm{dens}(R)        &= \frac{|\{\,e\in E : \text{\textenglish{both endpoints}} \in M\,\}|}{|M|}. \label{eq:dens}
\end{align}
الحرارة مقدار «يقظة» المنطقة؛ والكثافة تماسكها السببي الداخلي (الحواف الداخلية لكل عضو).
ويُدفِئ التفعيل منطقةً ($\mathrm{temp}\mathrel{+}=0.25$، بحدٍّ أقصى) ويُسجِّل نبضةً منطقية؛
وتضرب خطوةُ تبريدٍ الحرارات في عامل تلاشٍ وتُخلي المناطق دون عتبةٍ دنيا.

\subsection{سياسة القبول الديناميكية}\label{sec:policy}

تصير العتبة الثابتة $0.6$ دالّةً في ضغط الرموز $u\in[0,1]$ (الجزء المستخدَم من الميزانية)
وتعقيد المهمّة $\kappa\in[0,1]$:
\begin{align}
\theta(u,\kappa)       &= \mathrm{clamp}(0.6 + 0.3\,u - 0.2\,\kappa,\; 0.05,\; 0.95), \\
a(R)                   &= 0.55\,\mathrm{temp}(R) + 0.30\,\frac{\mathrm{dens}(R)}{1+\mathrm{dens}(R)} + 0.15\,\kappa, \\
\mathrm{admit}(R)      &\iff a(R) \ge \theta(u,\kappa).
\end{align}
يمكن قبول منطقةٍ ساخنةٍ متماسكة حتى حين يرفضها الثابت $0.6$؛ وترفع النافذةُ المُمتلئة تقريبًا
العتبةَ $\theta$ فلا يدخل سوى أسخن المناطق.

\section{الحتمية وإعادة التشغيل}\label{sec:determinism}

\begin{theorem}[الحتمية الإقليمية]
تجزئة المناطق، وكل قياس منطقةٍ في المعادلات~\eqref{eq:temp}--\eqref{eq:dens}، وتاريخ
التفعيل، كلّها دوالّ خالصة مستقلّة عن الترتيب في الرسم $G$ وفي متتالية أحداث التفعيل (ذات
الطوابع الزمنية المنطقية). وبالتالي، بمعلومية سجلّ أحداث جلسةٍ ما، تُعاد بناء حالة المحرّك
على نحوٍ مطابق: إن كان $G'$ هو الرسم المُعاد بناؤه من السجلّ و$L$ أحداثه الإقليمية، فإنّ
$\mathrm{replay\_from}(G',L)$ يساوي المحرّك الحيّ الذي أنتج $L$.
\end{theorem}
\begin{proof}[مخطّط البرهان]
يَعُدّ التجميع العُقد والحواف بترتيبٍ مُرتَّب ويحسب المركّبات المتّصلة ببحثٍ عرضي مُرتَّب،
فيكون خرجه مستقلًّا عن ترتيب تكرار التجزئة. والمعادلات \eqref{eq:score} و\eqref{eq:temp}
و\eqref{eq:dens} حسابٌ حتمي على حقول العُقد وعَدٍّ للحواف المؤهَّلة. ويقرأ التفعيل ساعةً
منطقية (عدّادًا)، لا ساعة الحائط أبدًا، ويُسجِّل حدثُ \code{RegionActivated} المُنبعِث
النبضةَ الدقيقة. وتُعيد البناءُ تجميعَ $G'$ (الحالة الأساس متطابقة، إذ $G'\!=\!G$ بنيويًّا)
ثم تُطبِّق التفعيلات المُسجَّلة بترتيب السجلّ؛ وكل خطوةٍ هي الدالّة الخالصة نفسها على المُدخلات
نفسها، فالحالة الناتجة متطابقة. ويتحقّق اختبار التكامل
\code{replay\_reconstructs\_identical\_engine} من $\mathrm{engine}=\mathrm{replay\_from}(G',L)$
بالمساواة البنيوية. \qed
\end{proof}

يوسّع هذا ضماناتِ CCOS القائمة: فسجلّ الأحداث الأساسي مُسلسَلٌ بالتجزئة ويكشف العبث، فتاريخ
المناطق قابلٌ للتدقيق، ولا يُظهِر اختبارٌ من 10\,000 دورة أيّ انحرافٍ في عدد المناطق أو
حراراتها.

\section{الجلسات القائمة على مصدر الأحداث والتنقيح بالسفر عبر الزمن}\label{sec:timetravel}

ليست الحتمية خاصّية صحّةٍ فحسب؛ بل هي أساس القدرة التي --- كما سيُحاجّ تقييمنا --- تميّز CCOS
عن مكدّس الاسترجاع. يُسجِّل \code{AgentSession} كل عمليةٍ إدراكية يجريها وكيلٌ على ذاكرته ---
\code{Ingest}، \code{SignalFailure}، \code{Recall} --- كخطٍّ زمني مُرتَّب. ولأنّ كل عملية
دالّةٌ حتمية في الحالة السابقة، تُعاد بناء الذاكرة بعد أوّل $n$ عملية بالضبط بإعادة تشغيل
العمليات المُحوِّرة منها (\code{replay\_to(n)}): إذ يمكن \emph{إرجاع عقل الوكيل} إلى أيّ
خطوة.

والعملية المهمّة للتنقيح هي الافتراض المضادّ. فـ\code{recall\_what\_if($n$, $q$, $b$)}
يُعيد تشغيل الذاكرة إلى الخطوة $n$ ويُعيد تنفيذ استدعاءٍ باستعلامٍ $q$ أو ميزانية رموز $b$
مختلفة، مُعيدًا النافذة التي \emph{كان} الوكيل سيراها. فحين يُصدِر وكيلٌ رقعةً سيّئة عند
الخطوة 15، يستطيع مُشغِّلٌ أن يُعيد تشغيل سياقه الدقيق عند الخطوة 14، ويوسّع الميزانية أو
يغيّر المرساة، ويُلاحِظ هل يتحسّن القرار --- مُنقِّحٌ مغلق الحلقة وقابلٌ للإعادة لذاكرة عمل
الوكيل. أمّا مكدّس RAG أو الإطار فيُحوِّر مخزنه احتماليًّا عبر الجلسة ولا يحفظ أيّ سجلّ
انتقالاتٍ معياري قابلٍ للإعادة، فالسؤال نفسه (\emph{لماذا، ومتى، فسد التمثيل؟}) لا جواب له
هناك. لا ندّعي أنّ هذا يُحسِّن نجاح المهام؛ بل ندّعي أنّه خاصّيةٌ لا تملكها البدائل، وأنّه
الموضع الأمين لإسهام CCOS (\S\ref{sec:protocol}).

\section{التنفيذ}\label{sec:impl}

المحرّك نحو $600$ سطرٍ من Rust خفيفة الاعتماديات، مُطبَّقةٍ فوق الرسم، من دون تغيير المُحلِّل
أو الحارس أو الباني التزايدي أو نواة مصدر الأحداث. والواجهة العامّة هي:
\code{ContextRegionEngine} (\code{cluster\_nodes}، \code{initialize\_regions}،
\code{activate\_region}، \code{tick\_cooldown}، \code{replay\_from})، وأنواع البيانات
\code{ContextRegion} / \code{ContextPoint} / \code{ContextWindow}، و\code{ContextPolicy}،
وخمسة متغيّرات أحداثٍ جديدة (\code{RegionCreated}، \code{RegionActivated}،
\code{RegionMerged}، \code{RegionEvicted}، \code{ContextWindowGenerated}). وتكشف واجهة
سطر الأوامر \code{ccos regions} التجميعَ والتفعيلَ وتقرير المحلّية. وتُبنى الحُزمة كاملةً من
دون تحذيراتٍ تحت \code{clippy -D warnings} وتجتاز 364 اختبارًا.

\section{التقييم: ما يمكننا قياسه اليوم}\label{sec:eval}

نفصل النتائج \emph{المقاسة} (هذا القسم، قابلٌ للإعادة كاملًا ومن دون نموذج لغوي) عن المكاسب
\emph{المفترَضة} على مستوى الوكيل (\S\ref{sec:protocol}).

\paragraph{الإعداد.} نُشغّل على شجرة مصدر CCOS نفسها ($|V|=705$ عقدة، $|E|=822$ حافة، تُنتِج
$26$ منطقة). ولعقدةٍ هدفٍ $t$ نأخذ $\Nk(t)$ بـ$k=2$ كحقيقةٍ مرجعية ونقارن استراتيجيتي
اختيارٍ عند ميزانية حجم المنطقة: \textbf{المسطّحة} --- أعلى $|R|$ عقدة بـ$\mathrm{score}$
(CCOS v0.2 / الاسترجاع المُرتَّب الكلاسيكي) --- و\textbf{المنطقة} --- أعضاء منطقة $t$. ونُبلِغ
عن الدقّة السببية $|S\cap\Nk|/|S|$، والاستدعاء $|S\cap\Nk|/|\Nk|$، والرموز التي تحتاجها كل
استراتيجية لـ\emph{تغطية} $\Nk$. وكل الأرقام يُصدِرها \code{scripts/region\_benchmark.sh}
وهي حتمية.

\begin{table}[h]
\centering
\begin{LTR}
\begin{tabular}{lcc}
\toprule
\textenglish{Strategy} & \textenglish{causal precision} & \textenglish{causal recall} \\
\midrule
\textenglish{flat (v0.2 / ranked)} & 0.021 & 0.347 \\
\textbf{\textenglish{region (v0.3)}} & \textbf{0.057} & \textbf{0.972} \\
\bottomrule
\end{tabular}
\end{LTR}
\caption{المحلّية على \code{src/} ($k=2$، 12 هدفًا). يغطّي الاختيار بالمناطق $97\%$ من
الجوار السببي لمهمّةٍ مقابل $35\%$ للمسطّح، وبرموزٍ أقل بنحو $48\%$ لتغطيةٍ مماثلة.}
\label{tab:locality}
\end{table}

\paragraph{التماسك والكلفة.} تبلغ المناطق \textbf{كثافةً سببية وسطية قدرها $0.955$}
(المعادلة~\ref{eq:dens}، بعد التطبيع) --- فهي متّصلةٌ داخليًّا بالكامل تقريبًا، أي عناقيد
سببية حقيقية لا تجميعات ملفاتٍ اعتباطية. ويستغرق بناء خريطة المناطق لـ$705$ عقدة نحو
$20$\,ms ($54.5$ عملية بناء/ثانية)؛ والتفعيل الواحد نحو $20$\,ms.

\paragraph{قراءةٌ أمينة.} الدقّة المطلقة منخفضة ($0.06$) لأنّ مُحلِّل v0.2 لا يُصدِر سوى
حواف الاحتواء والاستيراد، فـ$\Nk(t)$ صغيرٌ جدًّا (غالبًا $2$--$3$ عُقد) بينما تمتدّ المنطقة
على ملفٍّ كامل. والمكاسب المتينة الحقيقية هي \emph{الاستدعاء} ($0.97$ مقابل $0.35$)
و\emph{كفاءة الرموز} (نحو $48\%$): فالمنطقة تحتوي بموثوقية جوار المهمّة السببي وتدفع رموزًا
أقل لذلك. ومن شأن حوافَ دلالية أغنى (رسم الاستدعاءات / تدفّق البيانات، البند P1.3 من خارطة
الطريق) أن تُشحِذ $\Nk$ وهي الرافعة الأساس لرفع الدقّة؛ نُشير إلى هذا بدل إخفائه.

\section{محاكاة الفرضية تحت أوراكل مُعلَن (من دون نموذج لغوي)}\label{sec:sim}

الأطروحة الجوهرية --- \emph{هل تُعين الذاكرة السببية بالمناطق وكيلًا في المهام الطويلة
المتعدّدة الملفات؟} --- تحتاج في النهاية إلى عمليات نموذجٍ لغوي (\S\ref{sec:protocol}). لكنّ
\emph{شرطها الضروري} قابلٌ للاختبار الآن، من دون نموذج لغوي: \textbf{لا يستطيع وكيلٌ حلّ
مهمّةٍ يغيب سياقها السببي المطلوب عن نافذته}. نقيس، تحت أوراكل صريح، هل تضع كل استراتيجية
اختيارٍ \emph{السياق السببي المطلوب في النافذة}.

\paragraph{الإعداد.} نُولِّد مستودعاتٍ اصطناعية وحدوية: $S$ نظامًا فرعيًّا مستقلًّا، كلٌّ
مجموعةُ ملفاتٍ تربطها سلسلةٌ سببية واحدة عبر الملفات إضافةً إلى رموز \emph{شَرَك} عالية
الدرجة؛ ولا حواف بين الأنظمة الفرعية، فكل نظامٍ فرعي منطقةٌ سببية محدودة. والبنية السببية
\emph{مفصولة} عن البنية المعجمية --- فجارُ السلسلة لا يشارك الهدف أيّ رمز معرّفٍ، بل حافةً
فقط --- لنمذجة اعتماديةٍ ضرورية سببيًّا لكنّها غير مشابهة معجميًّا. وتتطلّب مهمّةٌ بنصف
\emph{قطر} $d$ النافذةَ $\pm d$ على طول سلسلتها (حتى $2d{+}1$ ملفًّا)؛ وتسع الميزانيةُ نحو
نظامٍ فرعي واحد، أقلّ بكثيرٍ من المستودع. و\textbf{الأوراكل} هو $R(t)\subseteq S$.

والأهمّ أنّ خطوط الاسترجاع تُحدِّد موضع الشيفرة من \emph{استعلامٍ} نصّي، بينما ذاكرةٌ بأسلوب
نظام التشغيل \emph{ترسو} على إشارة مساحة العمل (الملفّ النشِط، اختبارٌ فاشل). ونُنمذِج كليهما:
كل مهمّةٍ تحمل استعلامًا (كيس رموز) ومرساةً (عقدة الملفّ النشِط)، ونُشغّل سيناريوهين ---
\textbf{نظيف} (الاستعلام يشير إلى الهدف) و\textbf{مشوّش} (شَرَكٌ خادعٌ في نظامٍ فرعي غير ذي
صلة يتفوّق معجميًّا على الهدف). وتتقاسم ست استراتيجيات الميزانية: \code{rag-dense} (أعلى
$k$ معجميًّا)، و\code{rag-hybrid} (معجمي $+$ درجة سببية)، و\code{graphrag-1hop} (أفضل إصابة
$+$ قفزة)، و\code{graphrag-bfs} (توسّعٌ غير محدود من أفضل إصابة)، و\code{ccos-from-query}
(منطقة CCOS لأفضل إصابةٍ \emph{معجمية} --- وهو استئصال)، و\code{ccos-region} (منطقة CCOS
لـ\emph{المرساة}). مُبذَّرةٌ وحتمية؛ تُعاد بـ\code{ccos experiment}.

\begin{table}[h]
\centering
\begin{LTR}
\begin{tabular}{lcccccc}
\toprule
& \multicolumn{4}{c}{\textenglish{success at diameter $d$}} & \multicolumn{2}{c}{\textenglish{overall}} \\
\cmidrule(lr){2-5}\cmidrule(lr){6-7}
\textenglish{Strategy} & $d{=}1$ & $d{=}2$ & $d{=}3$ & $d{=}4$ & \textenglish{succ.} & \textenglish{cov.} \\
\midrule
\multicolumn{7}{l}{\emph{\textenglish{Clean query (points at the target)}}}\\
\code{rag-dense}      & 0.00 & 0.00 & 0.00 & 0.00 & 0.00 & 0.19 \\
\code{rag-hybrid}     & 1.00 & 0.00 & 0.00 & 0.00 & 0.23 & 0.65 \\
\code{graphrag-1hop}  & 1.00 & 0.00 & 0.00 & 0.00 & 0.23 & 0.58 \\
\code{graphrag-bfs}   & 1.00 & 1.00 & 1.00 & 1.00 & 1.00 & 1.00 \\
\code{ccos-from-query}& 1.00 & 1.00 & 1.00 & 1.00 & 1.00 & 1.00 \\
\code{ccos-region}    & 1.00 & 1.00 & 1.00 & 1.00 & 1.00 & 1.00 \\
\midrule
\multicolumn{7}{l}{\emph{\textenglish{Noisy query (a decoy out-scores the target lexically)}}}\\
\code{rag-dense}      & 0.00 & 0.00 & 0.00 & 0.00 & 0.00 & 0.19 \\
\code{rag-hybrid}     & 1.00 & 0.00 & 0.00 & 0.00 & 0.23 & 0.65 \\
\code{graphrag-1hop}  & 0.00 & 0.00 & 0.00 & 0.00 & 0.00 & 0.19 \\
\code{graphrag-bfs}   & 0.00 & 0.00 & 0.00 & 0.00 & 0.00 & 0.00 \\
\code{ccos-from-query}& 0.00 & 0.00 & 0.00 & 0.00 & 0.00 & 0.00 \\
\textbf{\code{ccos-region}} & \textbf{1.00} & \textbf{1.00} & \textbf{1.00} & \textbf{1.00} & \textbf{1.00} & \textbf{1.00} \\
\bottomrule
\end{tabular}
\end{LTR}
\caption{محاكاة الفرضية ($800$ مهمّة، البذرة $42$، ميزانيةٌ بنحو نظامٍ فرعي واحد). النجاح
$=$ المجموعة السببية المطلوبة $R(t)$ داخل النافذة. تحت استعلامٍ نظيف تتعادل كل طريقةٍ واعية
بالبنية؛ وتحت استعلامٍ مُضلِّل لا تنجو سوى المنطقة المرتكزة على مساحة العمل.}
\label{tab:sim}
\end{table}

\paragraph{النتائج (وحدودها الأمينة).}
\textbf{(1)} يفشل الاسترجاع المعجمي (\code{rag-dense}) في المهام السببية عبر الملفات ($0\%$؛
التغطية $0.19$): فالتشابه لا يُظهِر سياقًا ضروريًّا سببيًّا لكنّه غير مشابهٍ معجميًّا --- وهي
مقدّمة الفرضية. (انتصارات \code{rag-hybrid} عند $d{=}1$ تأتي من \emph{الدرجة السببية} التي
يستعيرها، لا من التشابه المعجمي، ولذلك تثبت تحت التشويش.)
\textbf{(2)} قيمة الاختيار الواعي بالبنية \emph{تنمو مع نصف القطر}: فالتوسّع بقفزةٍ واحدة لا
يحلّ سوى $d{=}1$، بينما الترحيل السببي الكامل (\code{graphrag-bfs}، \code{ccos-region}) يحلّ
كل $d$ --- وهو اتجاه H2.
\textbf{(3) التعادل في الحالة النظيفة.} تحت استعلامٍ نظيف، تبلغ \code{graphrag-bfs}
و\code{ccos-from-query} و\code{ccos-region} كلّها $1.00$: فالرافعة هي \emph{البنية} السببية،
\textbf{لا CCOS بذاته}.
\textbf{(4) الانفصال في الحالة المشوّشة.} تحت استعلامٍ \emph{مُضلِّل}، تنهار كل طريقةٍ تُحدِّد
موضع الشيفرة معجميًّا إلى $0\%$ --- بما في ذلك \code{graphrag-bfs} القوية \emph{و}استئصال
\code{ccos-from-query} (CCOS المُبذَّر على الاستعلام). ولا تنجو سوى \code{ccos-region}، التي
ترسو على إشارة مساحة العمل لا على الاستعلام، بـ$1.00$. ويعزل الاستئصالُ المُفارِق: إنّه
\emph{مصدر المرساة} (إشارةٌ بنيوية لمساحة العمل مقابل استعلامٍ معجمي)، لا آلية المناطق. وهذا
هو النظام الواقعي --- فوصف المهمّة نادرًا ما يُسمّي السبب البعيد الدقيق، وذاكرةٌ بأسلوب نظام
التشغيل تتعقّب مجموعة العمل النشِطة متينةٌ حيث يُضلَّل الاسترجاع وقت الاستعلام.
\textbf{(5) الافتراضات، مُعلَنة.} يَنسب هذا إلى CCOS مرساةً موثوقة (الملفّ النشِط / الاختبار
الفاشل)، تملكها ذاكرةٌ على مستوى نظام التشغيل دون خط استرجاعٍ صرف؛ ويفترض مستودعاتٍ
\emph{وحدوية} ذات مناطق قابلة للفصل (كثافة $0.955$ على شيفرةٍ حقيقية، \S\ref{sec:eval}) ---
فمنطقةٌ عملاقة متآلفة تتجاوز الميزانية تُطبِق الميزة (لوحِظ، وأُبلِغ عنه).
\textbf{(6)} وفي كل ذلك هذه \emph{محاكاةٌ تحت أوراكل مُعلَن}: تختبر الشرط \emph{الضروري}
(الاسترجاع)، لا الشرط \emph{الكافي} (التوليد) من \S\ref{sec:protocol}.

\section{تقييمٌ بنموذجٍ لغوي حقيقي: قياسٌ أوّل على الجهاز}\label{sec:protocol}

أرست المحاكاة (\S\ref{sec:sim}) الشرط الضروري. أمّا الشرط الكافي --- أن يَحلّ وكيلٌ لغوي
بعدئذٍ مهامّ أكثر بالذاكرة الإقليمية منه باسترجاع المقاطع --- فيتطلّب عمليات نموذج. نحن
\textbf{نُنفِّذ المنصّة} (\code{ccos eval}) ونُبلِغ عن \textbf{أوّل تشغيلٍ بنموذجٍ لغوي
حقيقي} إزاء نموذجٍ محلّي.

\paragraph{المنصّة المُنفَّذة.} كل مهمّة مشروعٌ صغير متعدّد الملفات يُرمِّز \emph{سلسلةً سببية
حسابية}: ثابتٌ أساس في ملفٍّ يُحوَّل عبر سلسلةٍ من دوالّ أحادية السطر في ملفاتٍ منفصلة،
والسؤال «أيّ عددٍ صحيح تُرجِعه الدالّة الأخيرة؟» لا يُجاب عنه \emph{إلّا} بقراءة السلسلة
كاملةً (متضمّنةً السبب البعيد). فالتقدير إذن مطابقةٌ تامّة على عددٍ صحيح --- موضوعيٌّ وقابلٌ
للأتمتة، من دون تنفيذ شيفرة. وتُجمِّع الاستراتيجيات الست نفسها من \S\ref{sec:sim} نافذة
السياق ضمن ميزانية رموز (مع فصل الاستعلام النظيف/المشوّش)، وتُرسَل النافذة إلى نموذج (أيّ
منفذٍ متوافق مع OpenAI أو Anthropic-Messages أو Ollama)، ونُسجِّل ثلاثة مقاييس لكل نصف
قطر: \textbf{نجاح المهمّة} (العدد الصحيح الصحيح)، و\textbf{تغطية الأوراكل} (السلسلة المطلوبة
$\subseteq$ النافذة --- مستقلّة عن النموذج)، و\textbf{هلوسة الرمز} (يُحيل الجواب إلى دالّةٍ
غائبة عن المشروع).

\paragraph{الإعداد.} نُشغّل \code{ccos eval} على الجهاز فوق NVIDIA Jetson AGX Thor إزاء
\code{qwen2.5:7b-instruct} المُقدَّم عبر Ollama (حرارة $0$)، بـ$20$ مهمّة، البذرة $7$،
ميزانية $600$ رمز، أنصاف أقطارٍ $1$--$4$، في النظامين النظيف والمشوّش. ويُبلِغ
الجدول~\ref{tab:real} عن التشغيل.

\begin{table}[h]
\centering
\begin{LTR}
\begin{tabular}{lcccccr}
\toprule
& \multicolumn{4}{c}{\textenglish{success at diameter $d$}} & & \\
\cmidrule(lr){2-5}
\textenglish{Strategy} & $d{=}1$ & $d{=}2$ & $d{=}3$ & $d{=}4$ & \textenglish{cov.} & \textenglish{tok.} \\
\midrule
\multicolumn{7}{l}{\emph{\textenglish{Clean query (names the target function)}}}\\
\code{rag-dense}      & 0.12 & 0.00 & 0.00 & 0.00 & 0.00 & 519 \\
\code{rag-hybrid}     & 0.00 & 0.00 & 0.00 & 0.00 & 0.00 & 521 \\
\code{graphrag-1hop}  & 1.00 & 0.00 & 0.00 & 0.00 & 0.40 & 270 \\
\code{graphrag-bfs}   & 1.00 & 0.00 & 0.17 & 0.00 & 1.00 & 402 \\
\code{ccos-from-query}& 1.00 & 0.00 & 0.17 & 0.00 & 1.00 & 402 \\
\code{ccos-region}    & 1.00 & 0.00 & 0.17 & 0.00 & 1.00 & 402 \\
\midrule
\multicolumn{7}{l}{\emph{\textenglish{Noisy query (a decoy out-scores the target lexically)}}}\\
\code{rag-dense}      & 0.00 & 0.00 & 0.00 & 0.00 & 0.00 & 531 \\
\code{rag-hybrid}     & 0.00 & 0.00 & 0.00 & 0.00 & 0.00 & 521 \\
\code{graphrag-1hop}  & 0.00 & 0.00 & 0.00 & 0.00 & 0.00 & 165 \\
\code{graphrag-bfs}   & 0.00 & 0.00 & 0.00 & 0.00 & 0.00 & 165 \\
\code{ccos-from-query}& 0.00 & 0.00 & 0.00 & 0.00 & 0.00 & 165 \\
\textbf{\code{ccos-region}} & \textbf{1.00} & \textbf{0.00} & \textbf{0.17} & \textbf{0.00} & \textbf{1.00} & \textbf{402} \\
\bottomrule
\end{tabular}
\end{LTR}
\caption{أوّل تشغيلٍ بنموذجٍ لغوي حقيقي: \code{qwen2.5:7b-instruct} عبر Ollama، على الجهاز
(NVIDIA Jetson AGX Thor)، $20$ مهمّة، البذرة $7$، ميزانية $600$ رمز. «\textenglish{cov.}»
هي تغطية الأوراكل المستقلّة عن النموذج؛ و«\textenglish{tok.}» وسط رموز الدخل. والنجاح جوابٌ
صحيحٌ بمطابقةٍ تامّة.}
\label{tab:real}
\end{table}

\paragraph{النتائج.}
\textbf{(1) التغطية تنتقل من المحاكاة إلى النصّ الحقيقي.} يُعيد عمود التغطية المستقلّ عن
النموذج إنتاجَ الجدول~\ref{tab:sim} على محتوى ملفاتٍ حقيقي: فالـ RAG المعجمي يغطّي $0$،
وتغطّي الطرائق الواعية بالبنية $1.00$ تحت استعلامٍ نظيف، وتحت التشويش \emph{وحدها}
\code{ccos-region} تُبقي $1.00$ --- الشرط الضروري، مُؤكَّدًا خارج المحاكي.
\textbf{(2) النجاح محدودٌ بالنموذج لا بالسياق.} حيث التغطية $1.00$، لا يحلّ نموذج $7$B سوى
$d{=}1$ ($1.00$) وجزءٍ من $d{=}3$ ($0.17$)، فائتًا $d{=}2,4$: فمع وجود السلسلة كاملةً في
النافذة، يكون الفشل المتبقّي \emph{استدلالًا حسابيًّا} --- سقفُ قدرةٍ للنموذج، لا فشلَ
اختيار. والأهمّ أنّ لا طريقة تنجح قطّ \emph{دون} تغطية (النجاح $\subseteq$ التغطية)، وهو
الادّعاء الذي بُنيت المنصّة لاختباره.
\textbf{(3) الانفصال المشوّش يصمد على نموذجٍ لغوي حقيقي.} تحت استعلامٍ مُضلِّل، تنهار كل
طريقةٍ مدفوعةٍ بالاستعلام إلى $0\%$ نجاح --- بما فيها \code{graphrag-bfs} واستئصال
\code{ccos-from-query} --- بينما \code{ccos-region}، المرتكزة على إشارة مساحة العمل، هي
\emph{الوحيدة} ذات النجاح غير الصفري ($d{=}1$ عند $1.00$). بل تبدو الطرائق المدفوعة
بالاستعلام \emph{أرخص} تحت التشويش ($165$ مقابل $402$ رمز) بالضبط لأنّها تسترجع بثقةٍ الملفّ
الأصغر الخاطئ؛ و\code{ccos-region} تُنفِق ميزانيتها على السلسلة الصحيحة سببيًّا.
\textbf{(4) حدودٌ أمينة.} $20$ مهمّة (نحو $5$ لكل نصف قطر) عيّنةٌ صغيرة، ونموذج $7$B يُرضِخ
الشرط الكافي إلى الأرض. ويُتوقَّع لنموذجٍ متقدّم (\code{deepseek-v4-pro}، قيد العمل) أو
نموذجٍ محلّي بحجم $70$B أن يرفع النجاح حيثما كانت التغطية $1.00$، \emph{شاحذًا} --- لا مُغيِّرًا
--- الانفصال، إذ تُثبِّت التغطيةُ السقفَ أصلًا.

\paragraph{نحو مقاييس خارجية.} تعزل مجموعة السلسلة الحسابية سؤال الاختيار؛ والاختبارات
الخارجية الحاسمة هي حلّ قضايا SWE-bench~\cite{jimenez2023swebench} (تجتاز رقعةٌ الاختبارات
المخفيّة) ومجموعةٌ مضبوطة من \emph{العلل المتعدّدة الملفات} حيث يقع موضع العطل وسببه ونصف
قطر انفجاره في ملفاتٍ متمايزة. وتتقاسم خطوط الأساس النموذجَ اللغوي الأساس وميزانية الرموز:
RAG الكلاسيكي~\cite{lewis2020rag} (أعلى $k$ بجيب تمام على المقاطع)، وGraphRAG~\cite{edge2024graphrag}
(ملخّصات المجتمعات على رسم شيفرة)، وMemGPT~\cite{packer2023memgpt} (ذاكرةٌ مُرحَّلة بأسلوب
نظام التشغيل)، ووكيل LangGraph~\cite{langgraph} بمخزنٍ شعاعي، ومناطق CCOS. والمقاييس: معدّل
الحلّ، ورموز الدخل حتى النجاح، ومعدّل هلوسة إسناد الاقتباس مُتحقَّقًا منه إزاء الرسم المرجعي.
وقد بدأنا هذا على تاريخٍ حقيقي: يُنقِّب \code{scripts/causal\_validation} عن تثبيتات الإصلاح،
ويستخرج شجرة ما قبل الإصلاح، ويحقن العطل، ويُقدِّر
$R_{\mathrm{cov}} = |F_{\mathrm{target}} \cap \mathrm{WorkingSet}_K| / |F_{\mathrm{target}}|$
--- وهو نسبة ملفات الإصلاح التي تستردّها مجموعة العمل المحدودة. وعلى تاريخ هذا المستودع، كشف
التشغيلُ الأوّل قيدًا وإصلاحَه، وكلاهما مقيس: فبانتشار الفشل في الاتجاه الأمامي فقط كانت
$R_{\mathrm{cov}}$ ثابتة عند $0.33$ (لم يُستردّ سوى ملفّ البذرة، لأنّ الملفات المُعدَّلة معًا
مُستورِدون \emph{أماميون} لا يُبلَغون إلّا عبر محاور الاعتماديات). وحلّ الاستيرادات داخل
الحُزمة إلى حواف ملفّ$\to$ملفّ ونشر الفشل \emph{في الاتّجاهين} يُصلِح هذا. وعبر ثلاث حُزَم
ناضجة --- \code{fd} و\code{bat} و\code{hyperfine}، $70$ تثبيت إصلاحٍ مُنقَّب --- يكون الأثر
متّسقًا: فعند ميزانيةٍ كافية ($K{\ge}50$) يرفع التغييران $R_{\mathrm{cov}}$ إلى $0.85$--$1.0$
(من $0.50$--$0.84$ في الاتجاه الأمامي وحده)، بينما \emph{يخفّ} إلى $0.19$--$0.28$ عند ميزانيةٍ
ضيّقة $K{=}20$. \textbf{لكن إزاء خط الأساس البديهي، ليس هذا انتصارًا.} فبتشغيل RAG معجمي
كلاسيكي (جيب تمام TF-IDF على نصّ الملفات، مُستعلَمًا بملفّ العطل) عند \emph{ميزانية الملفات
نفسها}، تكون $R_{\mathrm{cov}}$ لـ CCOS\,/\,RAG هي $0.92/0.94$ و$1.00/0.98$ و$0.87/0.92$ عند
$K{=}50$ وتتعادل عند $K{=}100$، بينما عند $K{=}20$ يتقدّم RAG بوضوح ($0.20/0.56$ على
\code{bat}، $0.20/0.73$ على \code{hyperfine}). فالاختيار السببي إذن \emph{لا ميزة تغطيةٍ
صافية} له على التشابه المعجمي هنا، وهو أسوأ عند ميزانيةٍ ضيّقة: ففي العلل الحقيقية تتشابه
ملفات الإصلاح معجميًّا، فيستردّها TF-IDF أيضًا --- ومقدّمة \S\ref{sec:sim} عن سببٍ غير مشابهٍ
معجميًّا لعَرَضِه لا تتكرّر على هذه المستودعات. والقراءة الأمينة أنّ الشرط \emph{الضروري}
يصمد للأغلبية الكبرى لكنّه \emph{ليس} خاصًّا بـ CCOS؛ وأيّ ميزةٍ حقيقية لا بدّ أن تأتي من
الأنظمة التي لا يختبرها هذا الإعداد --- استعلامٌ مُتدهوِر/غائب (\S\ref{sec:sim})، أو عَرَضٌ
(اختبارٌ فاشل) بعيدٌ معجميًّا عن سببه (يحتاج البذرة الحقيقية المدفوعة بالاختبار، لا أُحادية
أعلى درجة)، أو الشرط \emph{الكافي} (المرحلة 4). كما أنّ مساحات العمل متعدّدة الحُزَم، القابلة
للربط الآن، تبقى للقياس على نطاقٍ واسع.

\paragraph{الشرط الكافي (المرحلة 4): الحلّ يتعادل، والكفاءة لا.} نُشغّل نصف التوليد على
إصلاحاتٍ حقيقية أحادية الملف من \code{fd}: يُعطى وكيلٌ سياقًا متساوي الميزانية مبنيًّا
بطريقتين --- المنطقة السببية لـ CCOS مقابل RAG معجمي أعلى $k$ --- ويُطلَب منه إعادة كتابة
الملفّ المُعتلّ، ويُقدَّر بنجاح \code{cargo test}، مع إعادة محاولةٍ و\emph{المُصرِّف في
الحلقة} (عند الفشل، \emph{عطل صفحة سياقٍ} يُحلِّل الخطأ بطبقة الأثر، ويحقن ضغطًا على الملفات
المُعطِلة، ويُعيد التوجيه بنافذةٍ مُنعَّشة). فبنموذجٍ ضعيف $7$B ومحاولةٍ واحدة، يحلّ كلاهما
$0/15$؛ وبـ\code{qwen3-coder-30B} وثلاث محاولات ترفع الحلقةُ كليهما إلى $2/10$ ($20\%$) ---
فتغذية عطل الصفحة الراجعة، لا المُسترجِع، هي ما يفكّ الحلّ، ولا يُظهِر CCOS \textbf{أيّ ميزة
حلٍّ} على RAG، اتّساقًا مع نتيجة الاسترجاع. أمّا \emph{الكفاءة} فتفصلهما: فعند حلٍّ متساوٍ
تبلغ منطقة CCOS المحدودة ذاتيًّا وسطيًّا $776$ رمز سياق مقابل $5366$ للمُسترجِع المعجمي
المالئ للميزانية --- خفضٌ بمقدار $\mathbf{6.9\times}$. ويصمد الأمر نفسه على نطاقٍ واسع من دون
نموذج: فعبر $51$ سيناريو إصلاحٍ أحادي الملف من \code{fd} و\code{bat} و\code{hyperfine} يُجمِّع
CCOS $700$--$1600$ رمز سياق مقابل نحو $6000$ لـ RAG، خفضٌ بمقدار $\mathbf{4.1}$--$\mathbf{9.1\times}$.
هذا هو المحور الوحيد الذي يهيمن فيه CCOS: لا \emph{ما} يسترجعه (يُضاهيه خط أساس) بل \emph{كم
قليلًا} يحتاج، لأنّ المنطقة السببية تتوقّف عند مجموعة العمل بدل حشو أعلى $k$ حتى الميزانية.
ونُصرِّح بالتحفّظات: عيّنة \emph{الحلّ} صغيرة جدًّا ($n{=}10$، تمريرتان)، وخط الأساس يملأ
الميزانية بالبناء (وRAG بـ$k$ مضبوط بعناية سيكون أكثر تناثرًا أيضًا). والادّعاء القابل للدفاع
أضيق لكنّه حقيقي: CCOS \emph{يُعايِر نفسه ذاتيًّا} --- يحدّ نفسه عند المنطقة السببية من دون
$k$ أو ميزانيةٍ تُضبَط --- فلا يُبدِّد النافذة قطّ، وهو بالضبط مغزى الترحيل عند الطلب.

\paragraph{الفرضيات.}
\textbf{H1} (الكفاءة): يبلغ CCOS نجاح مهمّةٍ متساوٍ برموز دخلٍ أقل من RAG/GraphRAG، لأنّ منطقةً
تغطّي الجوار السببي باستدعاءٍ أعلى (الجدول~\ref{tab:locality}).
\textbf{H2} (النجاح بعيد المدى): في المهام المتعدّدة الملفات يتجاوز معدّل حلّ CCOS الـ RAG
بالمقاطع بهامشٍ يكبر مع نصف القطر السببي للمهمّة.
\textbf{H3} (الإسناد): يخفّض CCOS معدّل هلوسة الرموز، لأنّ النافذة المقبولة رسمٌ فرعي متّصل من
عُقدٍ \emph{حقيقية}.

\paragraph{تهديدات الصلاحية.} جودة المنطقة محدودةٌ بجودة الحواف (\S\ref{sec:eval})؛ فالنتائج
قد تخلط \emph{سياسة} المنطقة بالـ\emph{رسم} الكامن. ولذا يجب على الاستئصالات أن تُثبِّت الرسم
وتُغيِّر استراتيجية الاختيار وحدها، وأن تُبلِغ عن نصف القطر السببي لكل مهمّة لتُنسَب المكاسب
إلى الاتّساق لا إلى استرجاع نصٍّ أكثر. واستئصال \code{ccos-from-query} في الجدول~\ref{tab:real}
هو بالضبط هذا الضابط --- الرسم والميزانية نفساهما، والمرساة وحدها مُبدَّلة --- وهو ينهار مع
خطوط الأساس المعجمية تحت التشويش. ونتيجةٌ إيجابية على H1--H3 ستشكّل الإسهام البحثي؛ ونتيجةٌ
صفرية ستظلّ تُثبِت البنية التحتية الحتمية القابلة للتدقيق.

\section{القيود}

نحن صريحون عمدًا بشأن ما \emph{لم} يُبرهَن.
\textbf{(1) لا ميزة تغطيةٍ على RAG في العلل الحقيقية.} المقياس الرئيس على البيانات الحقيقية
($R_{\mathrm{cov}}$، \S\ref{sec:protocol}) يتعادل مع مُسترجِع معجمي بسيط من نوع TF-IDF عند
ميزانيةٍ كافية ويخسر أمامه عند ميزانيةٍ ضيّقة؛ وفي تثبيتات الإصلاح الحقيقية تتشابه الملفات
التي يمسّها إصلاحٌ معجميًّا، فمقدّمة المحاكاة عن سببٍ غير مشابهٍ معجميًّا لا تصمد. والأرقام
المطلقة القوية شرطٌ \emph{ضروري} يُحقِّقه خط أساسٍ أيضًا، لا انتصارَ CCOS.
(ويَعزل قياسٌ مُكمِّل على مصدر المحرّك نفسه علاقةً مختلفة مُعرَّفةً بنيويًّا: إذ لا يستردّ مُسترجِعٌ معجمي من نوع TF-IDF سوى $49\%$ (الاستدعاء@10) من اعتماديات الاستيراد المُستبانة بالـ AST --- وهي أزواجٌ لا تتقاسم مفرداتٍ في المعتاد، جيب تمام الاعتمادية $0.53$ مقابل $0.43$ عشوائيًّا --- فنقطةُ RAG العمياء حقيقية، لكنّها تقع خارج محور تماسُك الإصلاح الذي يقيسه المقياس الرئيس، حيث تتقاسم ملفاتُ الإصلاحِ مفرداتٍ بالفعل. وتستردّ الطبقةُ البنيوية تلك الحواف بالبناء؛ والـ AST، الافتراضيُّ الآن، هو ما يجعل ذلك الاسترداد كاملًا.)
\textbf{(2) الضروري $\neq$ الكافي.} لا نقيس قطّ هل \emph{تُعين} منطقةٌ وكيلًا على \emph{إصلاح}
علّة (المرحلة 4: رقعةٌ تجتاز الاختبارات المخفيّة) --- بل فقط هل الملفات ذات الصلة قابلة
للاسترجاع.
\textbf{(3) أُحادية البذرة.} تحقن المنصّة العطل في الملفّ المُعدَّل ذي أعلى درجة خروج، لا في
الاختبار الفاشل فعلًا، فلا تُمارِس النظامَ الذي يكون فيه عَرَضٌ بعيدًا معجميًّا عن سببه --- وهو
النظام الأكثر مواتاةً للاختيار السببي.
\textbf{(4) النطاق والإحصاء.} ثلاثة مستودعاتٍ أُحادية الحُزمة، $n\!\approx\!20$ لكلٍّ؛ تُبلَغ
الفروق بانحرافها المعياري لكنّها غير مُتحقَّق منها تقاطعيًّا.
\textbf{(5) المحرّك.} يعتمد المُحلِّل الآن افتراضيًّا على شجرة AST حقيقية من \code{syn} ---
قِيسَ أنّها أدقُّ بمقدار $36.5\%$ من جوار السطر السابق على مصدر المحرّك نفسه (ثُلثا استدعاء
الاستيراد مقابل استدعائه كاملًا)، ولم يُبقَ على الجوار إلّا احتياطيًّا لمدخلاتٍ غير Rust ---
وصُلِّبَ الاستيعاب لتُعاد بناؤه بِتةً بِتة عبر العمليات (استبانةُ استيرادٍ مُرتَّبة؛ ومركزيةٌ
ثابتةٌ بإزاء الترتيب). وتظلّ المناطق بحبيبية الملفّ ولا تندمج إلّا على حواف عبر الملفات
صريحة؛ وأوزان الدرجة مضبوطةٌ يدويًّا، وإن صارت متاحةً الآن مركزيةُ المتجه الذاتي ودورةُ حياة
العُقدة (\code{Stable}/\code{Working}/\code{Orphan}) كتحسيناتٍ مُعطَّلةٍ افتراضيًّا. لا شيء من
هذا يمسّ الخصائص \emph{المُبرهَنة} (التجزئة، الحتمية، إعادة التشغيل) ولا المحلّية المقاسة ---
لكنّ الحُجّة التجريبية لفائدةٍ لاحقة للوكيل، عند هذه النقطة، \emph{غير مُبرهَنة}.

\section{خاتمة}

انطلقنا لنُبيِّن أنّ تنظيم ذاكرة وكيلٍ بـ\emph{المناطق السببية} يسترجع السياق بعيد المدى
أفضل من استرجاع المقاطع، وأعطينا البناء تعريفًا دقيقًا وحتميًّا، وبرهنّا أنّه يُعاد بناؤه بِتةً
بِتة عند إعادة التشغيل. ثم اختبرنا ادّعاء الاسترجاع إزاء خط الأساس البديهي على بياناتٍ حقيقية
--- فلم يصمد: فعبر $70$ تثبيتًا حقيقيًّا لإصلاح علل يتعادل مُسترجِعٌ معجمي بسيط من نوع TF-IDF
مع الاختيار السببي ويتفوّق عليه عند ميزانيةٍ ضيّقة، ويخسر محورُ أثر الانهيار أمام «RAG على
رسالة الخطأ»، لأنّ ملفات الإصلاح ورسائل خطئه على شيفرةٍ حقيقية تتشارك المفردات. نُبلِغ بالنتيجة
السلبية بدل دفنها. والباقي ليس مُسترجِعًا أفضل بل صنفًا مختلفًا من الأشياء: ذاكرة عملٍ
\emph{حتمية، قابلة لإعادة التشغيل، قابلة للتدقيق} يمكن فيها إرجاع حالة سياق الوكيل الدقيقة
وإعادة تشغيلها بمعاملاتٍ مختلفة --- تنقيحٌ بالسفر عبر الزمن لا يستطيع مكدّس استرجاعٍ احتمالي
تقديمه. وهل تُؤتي قابلية التدقيق هذه، أو ميزة سياقٍ مُهيكَل وقت \emph{التوليد} (المرحلة 4)،
مكسبًا لاحقًا قابلًا للقياس، فأمرٌ يبقى مفتوحًا. نُصدِر المحرّك ومنصّة التحقّق الأمينة وهذه
الورقة كي يُتاح التحقّق --- لا مجرّد الادّعاء --- من التمييز بين ما هو مُبرهَن، وما هو مقيس،
وما هو مُفنَّد.

\begin{thebibliography}{99}
\bibitem{lewis2020rag} P.~Lewis et al. Retrieval-Augmented Generation for
Knowledge-Intensive NLP Tasks. \emph{NeurIPS}, 2020. arXiv:2005.11401.
\bibitem{asai2023selfrag} A.~Asai et al. Self-RAG: Learning to Retrieve, Generate
and Critique through Self-Reflection. \emph{ICLR}, 2024. arXiv:2310.11511.
\bibitem{edge2024graphrag} D.~Edge et al. From Local to Global: A Graph RAG
Approach to Query-Focused Summarization. 2024. arXiv:2404.16130.
\bibitem{yamaguchi2014cpg} F.~Yamaguchi et al. Modeling and Discovering
Vulnerabilities with Code Property Graphs. \emph{IEEE S\&P}, 2014.
\bibitem{packer2023memgpt} C.~Packer et al. MemGPT: Towards LLMs as Operating
Systems. 2023. arXiv:2310.08560.
\bibitem{park2023generative} J.~S.~Park et al. Generative Agents: Interactive
Simulacra of Human Behavior. \emph{UIST}, 2023. arXiv:2304.03442.
\bibitem{langgraph} LangChain. LangGraph: Building Stateful, Multi-Actor
Applications with LLMs. Software framework, 2023--2024.
\url{https://github.com/langchain-ai/langgraph}.
\bibitem{denning1968working} P.~J.~Denning. The Working Set Model for Program
Behavior. \emph{Communications of the ACM}, 11(5), 1968.
\bibitem{liu2023lost} N.~F.~Liu et al. Lost in the Middle: How Language Models Use
Long Contexts. \emph{TACL}, 2024. arXiv:2307.03172.
\bibitem{kwon2023paged} W.~Kwon et al. Efficient Memory Management for Large
Language Model Serving with PagedAttention. \emph{SOSP}, 2023. arXiv:2309.06180.
\bibitem{haber1991timestamp} S.~Haber and W.~S.~Stornetta. How to Time-Stamp a
Digital Document. \emph{Journal of Cryptology}, 3(2), 1991.
\bibitem{jimenez2023swebench} C.~E.~Jimenez et al. SWE-bench: Can Language Models
Resolve Real-World GitHub Issues? \emph{ICLR}, 2024. arXiv:2310.06770.
\end{thebibliography}

\end{document}
